% ==============================================================================
% CHAPITRE 5 : DISCUSSION ET PERSPECTIVES
% ==============================================================================
\chapter{Discussion et Perspectives}

\begin{infobox}[Objectif du chapitre]
Ce chapitre propose une analyse critique des résultats obtenus, identifie
les limites actuelles du système et trace les perspectives d'évolution
et d'amélioration futures.
\end{infobox}

% ──────────────────────────────────────────────────────────────────────────────
\section{Analyse Critique des Résultats}
% ──────────────────────────────────────────────────────────────────────────────

\subsection{Points forts du système}

\subsubsection{Performance remarquable sur un petit jeu de données}

Obtenir une précision de 79.41\% avec seulement 225 exemples d'entraînement
est un résultat significatif. Il illustre la puissance du transfer learning :
CamemBERT, ayant été pré-entraîné sur 138 Go de texte français, possède
déjà une compréhension profonde de la langue. Le fine-tuning n'a qu'à
adapter cette compréhension générale à la tâche spécifique de la
classification de sentiments étudiants.

À titre de comparaison, un modèle LSTM entraîné à partir de zéro sur le
même jeu de données obtiendrait typiquement une précision inférieure à 60\%,
car il n'aurait pas le bénéfice de ce vaste pré-entraînement.

\subsubsection{Architecture microservices modulaire}

La séparation claire entre le modèle IA, l'API, le dashboard et
l'application mobile permet :
\begin{itemize}
  \item De mettre à jour le modèle sans toucher aux interfaces utilisateur.
  \item D'ajouter de nouvelles interfaces (ex: application web) sans
        modifier le backend.
  \item De scaler chaque composant indépendamment selon la charge.
  \item De remplacer facilement une technologie par une autre (ex: MySQL
        au lieu de PostgreSQL) sans impact sur les autres couches.
\end{itemize}

\subsubsection{Accessibilité multi-canal}

Le système répond à deux profils d'utilisateurs distincts :
\begin{itemize}
  \item \textbf{L'étudiant}, qui accède via une application mobile intuitive,
        sans barrière technique.
  \item \textbf{L'enseignant}, qui dispose d'un tableau de bord analytique
        riche en visualisations.
\end{itemize}

\subsection{Limites identifiées}

\subsubsection{Taille du jeu de données}

La principale limitation du système est la taille du jeu de données
d'entraînement (225 exemples). Bien que les résultats soient encourageants,
un modèle entraîné sur 10 000 exemples ou plus serait nettement plus robuste,
notamment pour :
\begin{itemize}
  \item Les expressions idiomatiques et le langage familier étudiant.
  \item Les commentaires très courts (1-2 phrases) qui sont plus difficiles
        à classifier.
  \item Les nuances culturelles spécifiques à certains établissements.
\end{itemize}

\subsubsection{Absence d'authentification}

Le système actuel ne dispose pas d'un mécanisme d'authentification.
Tout utilisateur connaissant l'URL de l'API peut soumettre des commentaires
ou accéder aux statistiques. Dans un contexte de production réel, cela
représente une vulnérabilité importante.

\subsubsection{Déploiement local uniquement}

Le système est actuellement configuré pour fonctionner en local. Pour
un déploiement en production accessible à tous les étudiants d'un
établissement, il faudrait déployer l'API sur un serveur cloud (AWS,
GCP, Azure, ou VPS) et adapter l'URL dans l'application mobile.

\subsubsection{Pas de gestion des doublons}

Le système actuel ne détecte pas les commentaires identiques soumis
plusieurs fois, ce qui peut fausser les statistiques.

% ──────────────────────────────────────────────────────────────────────────────
\section{Impact Potentiel du Projet}
% ──────────────────────────────────────────────────────────────────────────────

Ce système, dans un contexte de déploiement réel, pourrait avoir un impact
significatif sur la qualité de l'enseignement supérieur :

\begin{table}[H]
\centering
\caption{Impact potentiel selon les parties prenantes}
\label{tab:impact}
\renewcommand{\arraystretch}{1.5}
\begin{tabular}{p{3cm} p{10cm}}
\toprule
\rowcolor{bleuprincipal!15}
\textbf{Partie prenante} & \textbf{Bénéfice concret} \\
\midrule
Étudiants & Retour d'expérience immédiat, sentiment d'être écouté, interaction plus engageante \\
\rowcolor{gray!8}
Enseignants & Identification rapide des cours mal compris, ajustement pédagogique en temps réel \\
Administration & Vue macro de la satisfaction par département, allocation optimale des ressources \\
\rowcolor{gray!8}
Établissement & Amélioration continue de la qualité, différenciation concurrentielle \\
\bottomrule
\end{tabular}
\end{table}

% ──────────────────────────────────────────────────────────────────────────────
\section{Perspectives d'Évolution}
% ──────────────────────────────────────────────────────────────────────────────

\subsection{Améliorations à court terme}

\begin{enumerate}
  \item \textbf{Agrandissement du jeu de données :} collecter 500 à 1000
        commentaires réels anonymisés d'étudiants pour re-fine-tuner le modèle
        et améliorer sa précision à 85\%+.

  \item \textbf{Authentification JWT :} implémenter JSON Web Tokens pour
        sécuriser l'API. Chaque étudiant s'authentifie et les commentaires
        sont liés à un compte anonymisé.

  \item \textbf{Interface web responsive :} ajouter une page web React ou
        Vue.js pour les étudiants ne disposant pas d'un smartphone Android.

  \item \textbf{Déploiement Docker :} conteneuriser le backend, la base de
        données et le dashboard dans des conteneurs Docker orchestrés par
        Docker Compose pour faciliter le déploiement et la portabilité.
\end{enumerate}

\subsection{Améliorations à moyen terme}

\begin{enumerate}
  \item \textbf{Analyse des aspects (Aspect-Based Sentiment Analysis) :}
        au lieu d'attribuer un seul sentiment à tout un commentaire,
        identifier les aspects mentionnés (``le contenu du cours'',
        ``la pédagogie du professeur'', ``le rythme'') et attribuer
        un sentiment à chacun.

  \item \textbf{Extraction de mots-clés :} identifier automatiquement
        les termes les plus fréquemment associés à chaque classe de sentiment
        (TF-IDF, YAKE ou KeyBERT) pour comprendre \textit{pourquoi} un
        cours est bien ou mal perçu.

  \item \textbf{Application iOS :} Flutter étant cross-platform, l'extension
        à iOS nécessite principalement une configuration Xcode et un compte
        développeur Apple.

  \item \textbf{Alertes automatiques :} envoyer un email ou une notification
        push à l'enseignant lorsque le taux de commentaires négatifs dépasse
        un seuil configurable.
\end{enumerate}

\subsection{Vision à long terme}

\begin{definitionbox}[Évolution vers un système de recommandation pédagogique]
À terme, le système pourrait évoluer vers une plateforme d'intelligence
pédagogique capable de :
\begin{itemize}
  \item Corréler les sentiments avec les résultats académiques (notes,
        taux d'abandon) pour identifier les cours à risque.
  \item Recommander automatiquement des ressources pédagogiques
        complémentaires aux étudiants dont les commentaires révèlent des
        difficultés de compréhension.
  \item Générer des rapports pédagogiques automatiques pour les conseils
        d'enseignement, alimentés par les analyses de sentiments agrégées.
  \item Intégrer des modèles multimodaux prenant en compte non seulement
        le texte mais aussi la fréquence de participation et les patterns
        temporels des commentaires.
\end{itemize}
\end{definitionbox}

\section*{Conclusion du Chapitre}

Ce chapitre a présenté une analyse honnête et nuancée du système développé.
Ses points forts résident dans l'utilisation d'une technologie de pointe
(CamemBERT), une architecture modulaire bien conçue et une accessibilité
multi-canal. Ses limites principales — liées à la taille du dataset et à
l'absence d'authentification — sont clairement identifiées et des solutions
concrètes ont été proposées. Ce travail constitue une base solide sur
laquelle des évolutions significatives pourront être construites.

\cleardoublepage
